\documentclass[aspectratio=169,10pt]{beamer}

% Shared Beamer setup for Fall 2026 course slides.
% Clean Cal Poly-inspired theme: no custom class files or uncommon packages.
\mode<presentation>{
  \usetheme{default}
  \usefonttheme{professionalfonts}
}

\definecolor{CalPolyGreen}{HTML}{154734}
\definecolor{CalPolyGold}{HTML}{BD8B13}
\definecolor{CalPolySage}{HTML}{7FA28F}
\definecolor{CalPolyMint}{HTML}{DDF1E7}
\definecolor{CalPolySand}{HTML}{A29F86}
\definecolor{CalPolyBeige}{HTML}{EFEEDE}
\definecolor{DeckInk}{HTML}{1F2933}
\definecolor{DeckMuted}{HTML}{5F6B7A}
\definecolor{DeckSoft}{HTML}{F7F7F5}

\setbeamersize{text margin left=0.55in,text margin right=0.55in}
\setbeamercolor{normal text}{fg=DeckInk,bg=white}
\setbeamercolor{structure}{fg=CalPolyGreen}
\setbeamercolor{title}{fg=white}
\setbeamercolor{frametitle}{fg=CalPolyGreen,bg=white}
\setbeamercolor{block title}{bg=CalPolySage,fg=white}
\setbeamercolor{block body}{bg=CalPolyMint,fg=black}
\setbeamercolor{alerted text}{fg=CalPolyGold}

\setbeamerfont{title}{size=\Huge,series=\bfseries}
\setbeamerfont{frametitle}{size=\Large,series=\bfseries}
\setbeamerfont{block title}{size=\normalsize,series=\bfseries}

\setbeamertemplate{navigation symbols}{}
\setbeamertemplate{itemize item}{\color{CalPolyGold}$\blacktriangleright$}
\setbeamertemplate{itemize subitem}{\color{CalPolyGreen}$\bullet$}
\setbeamertemplate{footline}{
  \hfill{\color{DeckMuted}\scriptsize\insertframenumber/\inserttotalframenumber}
  \hspace{0.18in}\vspace{0.12in}
}
\setbeamertemplate{frametitle}{
  \vspace{0.18in}
  \hspace{0.02in}\insertframetitle\par
  \vspace{0.08in}
  \noindent\makebox[\linewidth][l]{%
    {\color{CalPolyGold}\rule{0.55in}{2pt}}%
    {\color{CalPolyGreen}\rule{\dimexpr\linewidth-0.55in\relax}{2pt}}%
  }
}

\newcommand{\CourseNumber}{Course}
\newcommand{\CourseTitle}{Course Title}
\newcommand{\LectureNumber}{0}
\newcommand{\LectureTitle}{Lecture Title}
\newcommand{\LectureDate}{Date}
\newcommand{\CourseUnit}{Unit}

\newcommand{\coursenumber}[1]{\renewcommand{\CourseNumber}{#1}}
\newcommand{\coursetitle}[1]{\renewcommand{\CourseTitle}{#1}}
\newcommand{\lecturenumber}[1]{\renewcommand{\LectureNumber}{#1}}
\newcommand{\lecturetitle}[1]{\renewcommand{\LectureTitle}{#1}}
\newcommand{\lecturedate}[1]{\renewcommand{\LectureDate}{#1}}
\newcommand{\courseunit}[1]{\renewcommand{\CourseUnit}{#1}}

\author{Kirk Duran}
\institute{California Polytechnic State University, San Luis Obispo}
\date{\LectureDate}

\newcommand{\makedecktitle}{
  \begingroup
  \setbeamercolor{background canvas}{bg=CalPolyGreen}
  \begin{frame}[plain]
    \vfill
    \begin{center}
      {\usebeamerfont{title}\color{white}\LectureTitle\par}
      \vspace{0.18in}
      {\Large\color{white}Lecture \LectureNumber\par}
    \end{center}
    \vfill
    \noindent{\color{white}\rule{\linewidth}{1.2pt}}\par
    \vspace{0.08in}
    \noindent\colorbox{CalPolyGold}{\makebox[0.24\linewidth][c]{\color{white}\Large\CourseNumber}}
    \hspace{0.12in}{\color{white}\Large Kirk Duran}
  \end{frame}
  \endgroup
}

\newenvironment{keyidea}{\begin{block}{Key idea}}{\end{block}}
\newenvironment{activity}{\begin{block}{In class}}{\end{block}}
\newenvironment{cpdefinition}{\begin{block}{Definition}}{\end{block}}
\newenvironment{exampleblockcp}{
  \begingroup
  \setbeamercolor{block title}{bg=CalPolySand,fg=white}
  \setbeamercolor{block body}{bg=CalPolyBeige,fg=black}
  \begin{block}{Example}
}{
  \end{block}
  \endgroup
}

\newcommand{\imageplaceholder}[2][0.3in]{%
  \noindent\makebox[\linewidth][c]{%
    \fcolorbox{CalPolySage}{CalPolyBeige}{%
      \parbox[c][#1][c]{0.86\linewidth}{%
        \centering
        {\color{DeckMuted}\scriptsize Image placeholder: }%
        {\color{DeckInk}\scriptsize\ttfamily\detokenize{#2}}%
      }%
    }%
  }%
}

\newcommand{\sectionframe}[1]{
  \begingroup
  \setbeamercolor{background canvas}{bg=CalPolyGreen}
  \begin{frame}[plain]
    \vfill
    {\color{white}\LARGE\bfseries #1\par}
    \vspace{0.18in}
    {\color{CalPolyGold}\rule{0.22\paperwidth}{3pt}}
    {\color{white}\rule{0.62\paperwidth}{3pt}}
    \vfill
  \end{frame}
  \endgroup
}

\newcommand{\placeholdernote}{
  \begin{block}{Placeholder}
    Replace these bullets with the final lecture material, examples, and in-class activities.
  \end{block}
}

\usepackage[T1]{fontenc}
\usepackage{lmodern}
\usepackage{amsmath,amssymb}
\usepackage{graphicx}
\usepackage{booktabs}
\usepackage{array}
\usepackage{xcolor}
\pdfstringdefDisableCommands{\def\translate#1{#1}}
\graphicspath{{../assets/lecture-33/}}

% Cal Poly CS 4880, Lecture 33, November 16, 2026.
% Expanded after reading all EIGHT pages of the supplied Deep Q Learning lecture.
% Follow-up to Lecture 32 on tabular Q-learning. Approx. 50 minutes:
%  (1) motivation + DQN architecture 8 min; (2) stable DQN + real demo 15 min;
%  (3) Double DQN 8 min; (4) dueling + combinations 8 min;
%  (5) demonstrations, imitation, DAgger 9 min; (6) takeaways 2 min.
% All original illustrations generated locally; existing lecture-32 agent/policy
% recap figures explicitly reused. Reproducible toy model seeded at 4880.
% QWOP and Pokemon are retained as source-lecture motivating cases;
% do NOT infer their exact implementation from a video screenshot.
% Sources at end include original DQN, Double DQN, Dueling DQN, PER,
% DAgger and DQfD papers. No claims of Atari performance from the toy demo.

\newcommand{\lecturevisual}[2][2.15in]{%
  \par{\centering\includegraphics[width=.98\linewidth,height=#1,keepaspectratio]{#2}\par}%
}
\newcommand{\architecturestage}[1]{%
 {\centering\setlength{\fboxsep}{5pt}%
 \fcolorbox{CalPolyGreen}{DeckSoft}{\parbox{.86\linewidth}{\centering\small #1}}\par}%
}
\newcommand{\conceptlead}[1]{{\small\bfseries\color{CalPolyGreen}#1\par}}
\newenvironment{calculation}[1][Mathematical description]{%
  \par\begingroup\setbeamercolor{calculation body}{bg=DeckSoft,fg=DeckInk}%
  \begin{beamercolorbox}[sep=4pt,wd=\linewidth]{calculation body}%
  \footnotesize{\bfseries\color{DeckMuted}#1\par}%
  \setlength{\jot}{0pt}\setlength{\abovedisplayskip}{3pt}\setlength{\belowdisplayskip}{3pt}%
}{\end{beamercolorbox}\endgroup}

% Keep this dense source deck readable within the shared course theme.
\setbeamerfont{frametitle}{size=\large,series=\bfseries}
\setbeamertemplate{frametitle}{%
  \vspace{.08in}\hspace{.02in}\insertframetitle\par\vspace{.04in}%
  \noindent\makebox[\linewidth][l]{%
    {\color{CalPolyGold}\rule{.55in}{2pt}}%
    {\color{CalPolyGreen}\rule{\dimexpr\linewidth-.55in\relax}{2pt}}}%
}
\renewenvironment{keyidea}{%
  \par\vspace{3pt}\begin{beamercolorbox}[sep=5pt,wd=\linewidth]{block body}%
  {\bfseries\color{CalPolyGreen}Key idea.}\ %
}{\end{beamercolorbox}}
\setbeamersize{text margin left=.42in,text margin right=.42in}
\coursenumber{CS 4880}
\coursetitle{Artificial Intelligence}
\lecturenumber{33}
\lecturetitle{Deep Q-Learning: Double, Dueling, and Imitation Learning}
\newcommand{\LectureDisplayTitle}{Deep Q-Learning:\\[3pt]Double, Dueling,\\[3pt]and Imitation Learning}
\lecturedate{November 16, 2026}
\courseunit{Learning Agents and Machine Learning}

\begin{document}
\begingroup
\renewcommand{\LectureTitle}{\LectureDisplayTitle}
\makedecktitle
\endgroup

\section{From Q-Tables to Neural Networks}

\begin{frame}{Last Lecture: Q-Learning Makes Decisions From Experience}
\small
\conceptlead{We learned to estimate an action's long-term value, then take the action with the greatest estimate.}
\begin{columns}[T,onlytextwidth]
\begin{column}{.50\textwidth}
\lecturevisual[2.10in]{previous_policy.png}
\end{column}
\begin{column}{.47\textwidth}
\begin{calculation}[Tabular Q-learning]
\[\begin{aligned}y&=r+\gamma\max_{a'}Q(s',a'),\\
Q(s,a)&\leftarrow Q(s,a)+\alpha[y-Q(s,a)].\end{aligned}\]
\end{calculation}
\begin{itemize}
\item Q-values are stored in a separate row for every state.
\item $\varepsilon$-greedy balances trying actions and choosing the current best action.
\item We trained by interacting with a finite 5$\times$5 grid.
\end{itemize}
\end{column}
\end{columns}
\end{frame}

\begin{frame}{Why a Q-Table Does Not Scale to Images}
\small
\conceptlead{A screenshot is not a convenient index into a finite lookup table.}
\begin{columns}[T,onlytextwidth]
\begin{column}{.55\textwidth}
\lecturevisual[2.32in]{state_explosion.png}
\end{column}
\begin{column}{.42\textwidth}
\begin{itemize}
\item A 5$\times$5 grid has 25 positions.
\item Tracking 12 binary dirt bits yields up to $25\cdot 2^{12}$ distinct position/dirt states.
\item A single 84$\times$84 8-bit image has $256^{7056}$ possible raw images.
\item Similar screenshots should share what an agent has learned.
\end{itemize}
\end{column}
\end{columns}
\begin{keyidea}Replace explicit values for every state with a parameterized function that \emph{generalizes}.\end{keyidea}
\end{frame}

\begin{frame}{The Core Idea: Approximate $Q(s,a)$ With a Neural Network}
\small
\lecturevisual[2.13in]{table_to_network.png}
\begin{itemize}
\item A network $Q(s,a;\theta)$ predicts action values from a state representation.
\item Related states can share learned features, unlike independent lookup-table rows.
\item The \emph{objective is still Q-learning}; the representation and optimization method change.
\end{itemize}
\end{frame}

\begin{frame}{Why DeepMind Used Atari Games}
\small
\conceptlead{The original DQN program challenged agents to learn game control from screen pixels and reward.}
\begin{columns}[T,onlytextwidth]
\begin{column}{.56\textwidth}
\lecturevisual[2.65in]{breakout_schematic.png}
\end{column}
\begin{column}{.41\textwidth}
\begin{itemize}
\item Early work showed an agent could learn across many Atari games using a common deep Q-network architecture.
\item The task combines visual perception, delayed rewards, and discrete joystick actions.
\item The agent is not handed a rule such as ``move below the ball.''
\end{itemize}
\end{column}
\end{columns}
\begin{keyidea}DeepMind's DQN (Mnih et al., 2015) made high-dimensional sensory inputs usable for Q-learning.\end{keyidea}
\end{frame}

\begin{frame}{What Goes Into and Comes Out of a DQN?}
\small
\conceptlead{For a fixed discrete action set, one network evaluation yields a value for each available action.}
\lecturevisual[1.99in]{dqn_cnn_architecture.png}
\begin{keyidea}The \textbf{chosen action} is $\arg\max_a Q(s,a;\theta)$; the network outputs values, not a guaranteed correct move.\end{keyidea}
\end{frame}

\begin{frame}{A Screenshot May Not Tell Us Which Way the Ball Moves}
\small
\conceptlead{A useful observation must encode motion, not just appearance at one instant.}
\lecturevisual[1.96in]{frame_stack.png}
\begin{itemize}
\item Typical Atari DQN preprocesses frames and stacks four successive observations.
\item Convolutions learn spatial patterns; temporal stacking distinguishes moving objects.
\item This is an approximation to the underlying game state, not an assurance of perfect observability.
\end{itemize}
\end{frame}

\begin{frame}{Breakout: Why Long-Term Value Matters}
\small
\begin{columns}[T,onlytextwidth]
\begin{column}{.55\textwidth}
\lecturevisual[2.55in]{breakout_schematic.png}
\end{column}
\begin{column}{.42\textwidth}
\begin{itemize}
\item The paddle must anticipate where the ball will be, not just where it is now.
\item Breaking a brick can give a reward; preserving the rally makes future rewards possible.
\item Deliberately opening a channel through bricks may be better than maximizing the next hit.
\item The game score provides the training signal, not action-by-action human labels.
\end{itemize}
\end{column}
\end{columns}
\end{frame}

\section{Stabilizing Deep Q-Learning}

\begin{frame}{DQN Still Uses a Bellman Backup}
\small
\conceptlead{The network tries to predict a reward plus the best future value.}
\begin{calculation}[One-step DQN target; $d=1$ at termination]
\[y=r+\gamma(1-d)\max_{a'}Q(s',a';\theta^-).\]
\end{calculation}
\begin{calculation}[Predict one selected action; update online weights $\theta$]
\[\mathcal L(\theta)=\mathbb E_{(s,a,r,s',d)\sim\mathcal D}\left[\left(y-Q(s,a;\theta)\right)^2\right].\]
\end{calculation}
\begin{itemize}
\item $\theta$ is the network currently being optimized; $\theta^-$ is a separate target copy.
\item A terminal transition has \textbf{no future Q-value term}.
\item The target is treated as a constant when differentiating the online-network loss.
\end{itemize}
\end{frame}

\begin{frame}{Why Naive Neural Q-Learning Can Be Unstable}
\small
\conceptlead{Three problems meet when a network learns by interacting with its own evolving policy.}
\begin{enumerate}
\item \textbf{Correlated data:} consecutive frames and actions look similar; ordinary SGD assumptions are strained.
\item \textbf{Moving targets:} the model predicts a target partly computed by its own changing estimates.
\item \textbf{Bootstrapping under approximation:} erroneous predictions can feed later predictions.
\end{enumerate}
\begin{keyidea}DQN's key engineering choices are \textbf{experience replay} and a \textbf{target network}. They reduce instability; they do not guarantee convergence.\end{keyidea}
\end{frame}

\begin{frame}{Innovation 1: Experience Replay}
\small
\conceptlead{Instead of training on only the newest transition, remember and revisit earlier experience.}
\lecturevisual[2.15in]{replay_buffer.png}
\begin{itemize}
\item Store $(s,a,r,s',d)$ in a fixed-capacity replay memory $\mathcal D$.
\item Train on randomly sampled minibatches, weakening correlations and reusing data.
\end{itemize}
\end{frame}

\begin{frame}{A Replay Minibatch Is Not a New Episode}
\small
\lecturevisual[1.80in]{random_batch.png}
\begin{itemize}
\item Replay sampling selects previously stored transitions in an order unrelated to their original time index.
\item The Bellman target still uses each \emph{stored transition's} own next state $s'$.
\item Reusing experience improves sample efficiency but makes the training data distribution depend on collection history.
\end{itemize}
\end{frame}

\begin{frame}{Innovation 2: A Separate Target Network}
\small
\conceptlead{Keep one Q-network for gradient steps and another as a temporarily frozen reference.}
\lecturevisual[1.53in]{two_networks.png}
\begin{calculation}[Hard synchronization after every $C$ environment steps]
\[\theta^-\leftarrow\theta.\]
\end{calculation}
\begin{keyidea}The target network is a \emph{delayed copy}, not an independently trained expert.\end{keyidea}
\end{frame}

\begin{frame}{Why Delayed Targets Make Optimization Easier}
\small
\lecturevisual[2.10in]{target_hold.png}
\begin{itemize}
\item Between synchronizations, target predictions remain fixed for a given next state.
\item The online model can reduce prediction error against a temporarily stationary target.
\item Copy too frequently and the target moves too fast; copy too rarely and it becomes stale.
\end{itemize}
\end{frame}

\begin{frame}{A Complete DQN Gradient Update}
\small
\lecturevisual[2.13in]{dqn_training_step.png}
\begin{itemize}
\item Replay supplies the current and successor observations.
\item The online network predicts $Q(s,a)$; the target network computes the bootstrap target.
\item Backpropagation updates \textbf{only} the online weights; synchronize target weights on schedule.
\end{itemize}
\end{frame}

\begin{frame}[fragile]{Algorithm: Put Acting, Replay, and Learning Together}
\footnotesize
\conceptlead{The agent repeatedly alternates between generating experience and optimizing its Q estimates.}
\begin{verbatim}
online = QNetwork(); target = deepcopy(online)
replay = ReplayBuffer(capacity); total_steps = 0
for episode in range(num_episodes):
    state = env.reset(); terminated = False
    while not terminated:
        action = epsilon_greedy(online(state), epsilon)
        next_state, reward, terminated = env.step(action)
        replay.add(state, action, reward, next_state, terminated)
        total_steps += 1
        if len(replay) >= batch_size:
            s, a, r, ns, done = replay.sample(batch_size)
            y = r + gamma * (1 - done) * target(ns).max(1).values
            pred = online(s).gather(1, a[:, None]).squeeze(1)
            optimizer_step(loss(pred, y.detach()))
        if total_steps % target_interval == 0:
            target.load_state_dict(online.state_dict())
        state = next_state
\end{verbatim}
{\scriptsize\color{DeckMuted}Pseudocode: practical systems also need replay warmup, tensor validation, and episode-truncation handling.}
\end{frame}

\begin{frame}{Demo: Train a Small DQN on a New GridWorld}
\small
\conceptlead{To see the machinery without an expensive Atari training run, train a real PyTorch network on a fixed maze.}
\begin{columns}[T,onlytextwidth]
\begin{column}{.45\textwidth}
\begin{itemize}
\item 5$\times$5 grid, three fixed wall squares, and a goal in the top-right corner.
\item State: 25-dimensional one-hot position vector; four discrete actions.
\item Reward: $+10$ on reaching the goal; $-0.15$ per legal nonterminal step; $-0.35$ for a blocked move.
\item Train from random nonterminal start squares.
\end{itemize}
\end{column}
\begin{column}{.52\textwidth}
\lecturevisual[1.73in]{toy_dqn_policy.png}
\end{column}
\end{columns}
\begin{keyidea}This demo is a different reward design from Lecture 32's $+100/-1$ example. It demonstrates the algorithm, not Atari-scale performance.\end{keyidea}
\end{frame}

\begin{frame}[fragile]{Demo Code: Build an Action-Value Network}
\small
\conceptlead{A small multilayer perceptron suffices for a position vector; pixels would call for convolutional features.}
{\footnotesize\begin{verbatim}
import torch
from torch import nn

class QNet(nn.Module):
    def __init__(self):
        super().__init__()
        self.net = nn.Sequential(
            nn.Linear(25, 48), nn.ReLU(),
            nn.Linear(48, 48), nn.ReLU(),
            nn.Linear(48, 4))

    def forward(self, state_batch):
        return self.net(state_batch)  # [batch, 4]
\end{verbatim}}
\begin{keyidea}A different input encoder changes the network architecture, but not the Q-learning interpretation of its outputs.\end{keyidea}
\end{frame}

\begin{frame}{Measured Demo Results: Learning Across Episodes}
\small
\lecturevisual[2.90in]{toy_dqn_results.png}
\begin{keyidea}In one seeded run: \textbf{420 training episodes}, \textbf{4,425 environment steps}, and a final \textbf{21/21 greedy evaluation success} over legal nonterminal starting squares.\end{keyidea}
\end{frame}

\begin{frame}{Measured Demo: Inspect the Learned Policy}
\small
\begin{columns}[T,onlytextwidth]
\begin{column}{.52\textwidth}
\lecturevisual[2.88in]{toy_dqn_policy.png}
\end{column}
\begin{column}{.45\textwidth}
\begin{itemize}
\item The figure uses the network's learned \textbf{greedy} action for each reachable position.
\item The gold rollout starts at the bottom-left and reaches the top-right goal in \textbf{8 actions}.
\item This is a \emph{single random seed on a small fixed map}, not evidence that DQN always learns efficiently.
\item In real tasks, evaluate over many runs and diverse initial conditions.
\end{itemize}
\end{column}
\end{columns}
\end{frame}

\section{Double Deep Q-Learning}

\begin{frame}{Problem: Taking a Maximum Can Overestimate Values}
\small
\conceptlead{Suppose every action is equally good, but our estimates contain random noise.}
\begin{calculation}[Even zero-mean errors can create positive maximization bias]
\[\mathbb E\!\left[\max_a\widehat Q(s',a)\right]\;\geq\;\max_a\mathbb E[\widehat Q(s',a)].\]
\end{calculation}
\begin{itemize}
\item The maximum tends to select an \emph{action whose estimation error happened to be positive}.
\item Standard DQN uses the same target-network values to \textbf{choose} the best next action and \textbf{value} it.
\item This can inflate bootstrap targets and sometimes damage action choice (van Hasselt et al., 2016).
\end{itemize}
\begin{keyidea}The problem is not ordinary training speed; it is an estimation bias introduced by the max operator.\end{keyidea}
\end{frame}

\begin{frame}{Worked Example: Selection and Evaluation Disagree}
\small
\lecturevisual[1.40in]{double_numeric_example.png}
\begin{calculation}[With $r=1$ and $\gamma=0.9$]
\[\underbrace{y_{\rm DQN}}_{\rm target\ max}=1+0.9(9.5)=9.55,\qquad
\underbrace{y_{\rm Double}}_{\rm online\ select}=1+0.9(6.0)=6.40.\]
\end{calculation}
\begin{keyidea}The \textbf{online network picks stay} (10.0); the \textbf{target network evaluates stay} (6.0), even though its own largest value is 9.5 for right.\end{keyidea}
\end{frame}

\begin{frame}{Double DQN: Decouple Choosing From Evaluating}
\small
\conceptlead{Use the online network for $\arg\max$, then the target network to evaluate that selected action.}
\begin{calculation}[Ordinary DQN target]
\[y_{\rm DQN}=r+\gamma(1-d)\max_{a'}Q(s',a';\theta^-).\]
\end{calculation}
\begin{calculation}[Double DQN target]
\[a^*=\arg\max_{a'}Q(s',a';\theta),\qquad
 y_{\rm DDQN}=r+\gamma(1-d)Q(s',a^*;\theta^-).\]
\end{calculation}
\begin{itemize}
\item \textbf{No additional independently trained network} is required beyond online and target.
\item The action for the backup comes from online estimates; its value comes from the delayed copy.
\item Double DQN \emph{reduces} overestimation; it does not make every value estimate unbiased.
\end{itemize}
\end{frame}

\begin{frame}{Measured Simulation: Why Independent Evaluation Helps}
\small
\lecturevisual[1.98in]{max_bias_montecarlo.png}
\begin{keyidea}This is a seeded Monte Carlo thought experiment with \textbf{20,000 trials}: all six actions truly have value zero and the noise is independent. It isolates maximization bias, \emph{not} learned DQN performance.\end{keyidea}
\end{frame}

\begin{frame}[fragile]{Code: Change Only How the Target Is Computed}
\small
\conceptlead{DQN and Double DQN can share replay memory, optimizer, state encoder, and training loop.}
{\footnotesize\begin{verbatim}
# batch tensors: states, actions, rewards, next_states, done
prediction = online(states).gather(1, actions[:, None]).squeeze(1)

with torch.no_grad():
    # Standard DQN:
    standard_next = target(next_states).max(dim=1).values

    # Double DQN: online selects; target evaluates
    best_next_actions = online(next_states).argmax(dim=1)
    double_next = target(next_states).gather(
        1, best_next_actions[:, None]).squeeze(1)

    y_double = rewards + gamma * (1.0 - done) * double_next
loss = torch.nn.functional.smooth_l1_loss(prediction, y_double)
\end{verbatim}}
\begin{keyidea}The \texttt{no\_grad} boundary prevents the target computation from becoming a gradient path.\end{keyidea}
\end{frame}

\section{Dueling Networks: Value and Advantage}

\begin{frame}{A Different Question: Is the State Good, or the Action Good?}
\small
\conceptlead{Imagine a racing car on a straight road: it may be in a good state even if most steering choices are nearly equivalent.}
\begin{columns}[T,onlytextwidth]
\begin{column}{.55\textwidth}
\begin{itemize}
\item \textbf{State value} $V(s)$ summarizes how promising the situation is.
\item \textbf{Action advantage} $A(s,a)$ measures how one action compares with the state's baseline.
\item Standard DQN predicts a Q-value for each action directly, without explicitly separating these concepts.
\item A dueling network gives the model separate representation streams for them.
\end{itemize}
\end{column}
\begin{column}{.42\textwidth}
\lecturevisual[2.19in]{dueling_architecture.png}
\end{column}
\end{columns}
\end{frame}

\begin{frame}{Dueling DQN: Shared Features, Two Heads}
\small
\lecturevisual[2.58in]{dueling_architecture.png}
\begin{keyidea}\textbf{Dueling describes the architecture}, not a new exploration rule or a replacement for Q-learning. The two streams recombine into ordinary Q-values.\end{keyidea}
\end{frame}

\begin{frame}{How Do Value and Advantage Combine?}
\small
\conceptlead{Naively $Q=V+A$ is not identifiable: a constant can move between the two terms.}
\begin{calculation}[Common mean-centered dueling aggregation]
\[Q(s,a)=V(s)+A(s,a)-\frac{1}{|\mathcal A|}\sum_{b\in\mathcal A}A(s,b).\]
\end{calculation}
\begin{itemize}
\item $V(s)$ is a scalar; $A(s,\cdot)$ has one output for each action.
\item Centering makes $V(s)$ equal to the mean predicted Q-value across actions.
\item Dueling is particularly helpful where many actions have similar consequences.
\end{itemize}
\end{frame}

\begin{frame}{Worked Example: One Good State, Different Actions}
\small
\lecturevisual[1.47in]{dueling_numeric.png}
\begin{calculation}[Illustration: V=5 and the mean advantage is 1/3]
\[Q=5+A-\overline A=\left[\tfrac{11}{3},\tfrac{14}{3},\tfrac{20}{3}\right]\approx[3.67,4.67,6.67].\]
\end{calculation}
\begin{keyidea}The right action is best, but the value head separately represents that this is an overall promising state.\end{keyidea}
\end{frame}

\begin{frame}[fragile]{Code: A Minimal Dueling Network in PyTorch}
\footnotesize
\conceptlead{Implement two output heads and center the action advantages along the action dimension.}
\begin{verbatim}
class DuelingQNet(nn.Module):
    def __init__(self, obs_dim=25, n_actions=4):
        super().__init__()
        self.features = nn.Sequential(
            nn.Linear(obs_dim, 64), nn.ReLU())
        self.value = nn.Linear(64, 1)
        self.advantage = nn.Linear(64, n_actions)

    def forward(self, s):
        h = self.features(s)
        v = self.value(h)                   # [batch, 1]
        a = self.advantage(h)               # [batch, actions]
        return v + a - a.mean(dim=1, keepdim=True)
\end{verbatim}
\begin{keyidea}The returned tensor still has one Q-value per action, so existing $\varepsilon$-greedy and replay code can be reused.\end{keyidea}
\end{frame}

\begin{frame}{Double + Dueling = Dueling Double DQN}
\small
\lecturevisual[2.10in]{variant_matrix.png}
\begin{itemize}
\item \textbf{Double DQN} changes the \emph{target calculation}: choose with online, value with target.
\item \textbf{Dueling DQN} changes the \emph{network architecture}: share features, split into value and advantage.
\item \textbf{Dueling Double DQN} uses both: two dueling network copies and the Double backup.
\end{itemize}
\end{frame}

\begin{frame}{Bonus: Prioritized Experience Replay}
\small
\conceptlead{Some experiences are more surprising than others; can training focus on them?}
\begin{columns}[T,onlytextwidth]
\begin{column}{.57\textwidth}
\lecturevisual[2.31in]{prioritized_replay.png}
\end{column}
\begin{column}{.40\textwidth}
\begin{itemize}
\item A common priority signal is absolute TD error $|y-Q(s,a)|$.
\item More surprising transitions are sampled more often, \emph{not exclusively}.
\item Importance-sampling weights can reduce sampling bias.
\item Replay priority is a separate design choice from Double and dueling.
\end{itemize}
\end{column}
\end{columns}
\begin{keyidea}``Prioritized Dueling Double DQN'' is a combination of independently motivated enhancements (Schaul et al., 2016).\end{keyidea}
\end{frame}

\section{Case Studies From the Original Slides}

\begin{frame}{Case Study 1: Breakout and the Pixel-to-Action Story}
\small
\begin{columns}[T,onlytextwidth]
\begin{column}{.53\textwidth}
\lecturevisual[2.42in]{breakout_schematic.png}
\end{column}
\begin{column}{.44\textwidth}
\begin{itemize}
\item The training signal is game score or other environment rewards.
\item Pixels encode paddle, ball, and bricks; four-frame input conveys velocity.
\item Delayed reward makes it rational to maneuver toward future brick-breaking opportunities.
\item This is the classic \emph{DQN} case study in the supplied deck.
\end{itemize}
\end{column}
\end{columns}
{\scriptsize\color{DeckMuted}Research: Mnih et al. (Nature, 2015). Illustration recreated for this lecture; not a gameplay capture.}
\end{frame}

\begin{frame}{Case Study 2: QWOP and the Difficulty of Control}
\small
\begin{columns}[T,onlytextwidth]
\begin{column}{.53\textwidth}
\lecturevisual[2.60in]{qwop_coordination.png}
\end{column}
\begin{column}{.44\textwidth}
\begin{itemize}
\item Each key controls part of a runner's legs; actions interact through physics.
\item A simple forward-progress reward may still require learning a delicate movement sequence.
\item The original slides use QWOP to motivate \textbf{Double DQN} and \textbf{prioritized replay}.
\item Exact algorithm details for the linked video are \emph{not independently established} here.
\end{itemize}
\end{column}
\end{columns}
\begin{keyidea}A neat demo is not enough to identify an implementation or to evaluate learning reliability.\end{keyidea}
\end{frame}

\begin{frame}{Case Study 3: Exploring a Long Game World}
\small
\begin{columns}[T,onlytextwidth]
\begin{column}{.53\textwidth}
\lecturevisual[2.55in]{pokemon_reward_map.png}
\end{column}
\begin{column}{.44\textwidth}
\begin{itemize}
\item The supplied lecture describes a Pokemon-playing agent receiving pixels and Game Boy-style actions.
\item Rewards may incentivize experience points, exploration, and battle progress.
\item Long-horizon games create sparse feedback and opportunities for reward hacking.
\item The cited creator's video is a \emph{case study}; we should not assume it used a standard DQN.
\end{itemize}
\end{column}
\end{columns}
{\scriptsize\color{DeckMuted}Case link from original slides: \href{https://www.youtube.com/watch?v=DcYLT37ImBY}{Peter Whidden, Training AI to Play Pokemon with Reinforcement Learning (2023)}.}
\end{frame}

\section{Imitation Learning: Learn From an Expert}

\begin{frame}{Why Not Learn Every Action From Scratch?}
\small
\conceptlead{Pure RL may spend enormous time on poor, unsafe, or uninformative exploration.}
\begin{itemize}
\item Many tasks have recordings of a human or an existing controller performing useful behavior.
\item \textbf{Imitation learning} exploits expert examples instead of initially discovering all good behavior from scalar reward.
\item A good demonstration can reduce early exploration, but it may not cover all states the agent will visit.
\item Human demonstrations are costly, may be inconsistent, and may not be optimal.
\end{itemize}
\begin{keyidea}The central additional information is a \textbf{demonstrated action} $a^*$ for an observation $s$.\end{keyidea}
\end{frame}

\begin{frame}{Behavior Cloning Is Supervised Learning of a Policy}
\small
\conceptlead{Train a network to predict the demonstrator's action from the observation.}
\lecturevisual[1.12in]{behavior_cloning_flow.png}
\begin{calculation}[Cross-entropy for a discrete expert action]
\[\mathcal L_{\rm BC}(\theta)=-\frac1N\sum_{i=1}^{N}\log\pi_\theta(a_i^*\mid s_i).\]
\end{calculation}
\begin{keyidea}Behavior cloning uses \emph{state--action labels}. It does not need rewards, a Q-network, or temporal-difference targets.\end{keyidea}
\end{frame}

\begin{frame}[fragile]{Code: Imitate Demonstrated Actions}
\small
\conceptlead{Imitation is a familiar classification problem; control becomes the test at deployment time.}
{\footnotesize\begin{verbatim}
# expert_states: [N, obs_dim]
# expert_actions: integer labels [N]
policy = nn.Sequential(
    nn.Linear(obs_dim, 64), nn.ReLU(),
    nn.Linear(64, n_actions))
optimizer = torch.optim.Adam(policy.parameters(), lr=1e-3)

for states, actions in expert_loader:
    action_logits = policy(states)
    loss = nn.functional.cross_entropy(action_logits, actions)
    optimizer.zero_grad()
    loss.backward()
    optimizer.step()

# At deployment, choose argmax of logits for a deterministic policy.
\end{verbatim}}
\begin{keyidea}This code learns to \emph{match labels}; only interaction testing reveals whether it recovers gracefully from mistakes.\end{keyidea}
\end{frame}

\begin{frame}{The Big Problem: Small Errors Change Future Observations}
\small
\lecturevisual[1.53in]{covariate_shift.png}
\begin{itemize}
\item The expert usually demonstrates near-successful trajectories.
\item One mistaken action can put the learner in a state absent from its training dataset.
\item New mistakes compound as the learner visits further unfamiliar states.
\end{itemize}
\begin{keyidea}High held-out action-prediction accuracy does not guarantee high closed-loop task success.\end{keyidea}
\end{frame}

\begin{frame}{DAgger: Have the Expert Label the Learner's States}
\small
\conceptlead{Dataset Aggregation addresses the distribution shift by collecting labels where the learner actually goes.}
\lecturevisual[1.75in]{dagger_loop.png}
\begin{itemize}
\item Roll out the current learner; ask the expert what it should have done at visited states.
\item Add these examples to the dataset, retrain, and repeat (Ross et al., 2011).
\item This can recover from unfamiliar situations, but requires an accessible expert or oracle.
\end{itemize}
\end{frame}

\begin{frame}{Combine Demonstrations With Reinforcement Learning}
\small
\lecturevisual[1.28in]{imitation_to_rl.png}
\begin{itemize}
\item \textbf{Option A:} pretrain a policy with behavior cloning, then fine-tune it through environment interaction.
\item \textbf{Option B:} seed or retain demonstrations in an off-policy replay buffer while learning Q-values.
\item \textbf{DQfD:} a documented DQN-with-demonstrations method combines supervised demonstration loss with TD learning, replay prioritization, and additional training choices.
\end{itemize}
\begin{keyidea}A classifier policy is not automatically a DQN: the architecture and loss must be connected deliberately when fine-tuning.\end{keyidea}
\end{frame}

\begin{frame}{Distinguish Imitation, Replay Seeding, and Reward Shaping}
\footnotesize
\begin{tabular}{@{}p{.23\linewidth}p{.33\linewidth}>{\raggedright\arraybackslash}p{.36\linewidth}@{}}
\textbf{Technique} & \textbf{Information provided} & \textbf{What actually changes}\\[3pt]\hline
Behavior cloning & Expert's action $a^*$ at state $s$ & Supervised policy objective\\[4pt]
DAgger & Expert labels on learner states & Dataset matches learner's visits better\\[4pt]
Demo replay / DQfD & Past expert transitions & Off-policy samples; sometimes extra losses\\[4pt]
Reward shaping & Intermediate numeric rewards & The reward function and learning signal\\[4pt]
\end{tabular}
\vspace{7pt}
\begin{keyidea}Reward shaping is \textbf{not itself imitation learning}; it supplies designed rewards rather than demonstrated expert actions.\end{keyidea}
\end{frame}

\section{Synthesis and Evaluation}

\begin{frame}{A Map of the Methods: What Does Each Improvement Fix?}
\footnotesize
\begin{tabular}{@{}p{.22\linewidth}p{.35\linewidth}p{.33\linewidth}@{}}
\textbf{Method} & \textbf{Key change} & \textbf{Remaining limitation}\\[3pt]\hline
Tabular Q-learning & State/action lookup table & Cannot scale to large inputs\\[4pt]
DQN & Neural Q + replay + target copy & Noisy targets and sample cost\\[4pt]
Double DQN & Online selects, target evaluates & Function approximation can still err\\[4pt]
Dueling DQN & Value and advantage streams & Still needs off-policy stabilization\\[4pt]
Dueling Double DQN & Both target and architecture changes & Hyperparameters and data matter\\[4pt]
Prioritized replay & Sample high-TD-error transitions more & Can bias sampling\\[4pt]
Behavior cloning & Supervised expert action labels & Compounding distribution shift\\[4pt]
DAgger / demos + RL & Learner-state labels / mixed training & Expert cost and integration complexity\\[4pt]
\end{tabular}
\end{frame}

\begin{frame}{Evaluate the Policy, Not Just the Training Loss}
\small
\lecturevisual[1.87in]{evaluation_protocol.png}
\begin{itemize}
\item Record episodic return, task completion, environment steps, and variability across random seeds.
\item Evaluate the greedy policy without training-time exploration; keep test cases held out.
\item For imitation, inspect both expert-action accuracy \emph{and closed-loop recovery}.
\item A loss curve that decreases is not proof of better real-world decisions.
\end{itemize}
\end{frame}

\begin{frame}{Check Your Understanding}
\small
\begin{enumerate}
\item If a single screenshot cannot reveal ball velocity, what information should the DQN receive?
\item Which network chooses the next action in \textbf{Double DQN}, and which network evaluates that action?
\item In a \textbf{dueling} network, why do we subtract the mean advantage?
\item A behavior-cloned robot makes one bad turn and enters an unfamiliar hallway. Which imitation-learning approach was designed to help with this?
\end{enumerate}
\vspace{9pt}
\begin{keyidea}The central architecture story: \textbf{approximate $Q$}, \textbf{stabilize backups}, \textbf{separate selection from evaluation}, \textbf{separate value from advantage}, and \textbf{learn from examples when available}.\end{keyidea}
\end{frame}

\begin{frame}{References, Videos, and Reproducibility}
\scriptsize
\begin{itemize}
\item Mnih et al. (2015). \emph{Human-level control through deep reinforcement learning}. \href{https://doi.org/10.1038/nature14236}{Nature 518:529--533}. Original DQN, replay, and frozen target network.
\item van Hasselt, Guez, Silver (2016). \emph{Deep Reinforcement Learning with Double Q-Learning}. \href{https://doi.org/10.1609/aaai.v30i1.10295}{AAAI 2016}.
\item Wang et al. (2016). \emph{Dueling Network Architectures for Deep Reinforcement Learning}. \href{https://proceedings.mlr.press/v48/wangf16.html}{ICML 2016}.
\item Schaul et al. (2016). \emph{Prioritized Experience Replay}. \href{https://arxiv.org/abs/1511.05952}{ICLR 2016}.
\item Ross, Gordon, Bagnell (2011). \emph{A Reduction of Imitation Learning and Structured Prediction to No-Regret Online Learning}. \href{https://proceedings.mlr.press/v15/ross11a.html}{AISTATS 2011}.
\item Hester et al. (2018). \emph{Deep Q-learning from Demonstrations}. \href{https://doi.org/10.1609/aaai.v32i1.11757}{AAAI 2018}.
\item Cases from the supplied source lecture: \href{https://www.youtube.com/watch?v=V1eYniJ0Rnk}{Breakout clip}; \href{https://www.youtube.com/watch?v=82sTpO_EpEc}{QWOP clip}; \href{https://www.youtube.com/watch?v=DcYLT37ImBY}{Pokemon clip}. Implementation claims for independently created videos are not verified by this deck.
\item Diagrams are locally generated schematics, with recap figures reused from Lecture 32. \texttt{generate-figures.py} produces visuals; \texttt{train-toy-dqn.py} generates the two measured toy learning figures. Seed 4880; no external datasets. Figures tagged \emph{illustrative} are not learned benchmarks.
\end{itemize}
\end{frame}

\end{document}
